%% RT2 Final Paper — Tactile Material Classification from Deformation & Shear
%% IEEE conference format (IEEEtran). Compile with pdflatex (x2) + bibtex if desired.
%%
%% NOTE TO AUTHOR (Amr): This paper is built on the RT2 method taught by
%% Prof. Recchiuto (PICO research question; independent/dependent/control
%% variables; careful baseline; multi-metric evaluation; figures with error
%% bars). Verify section requirements against the actual RT2 rubric before
%% submission. All numbers are the HONEST values from the project; do NOT
%% replace them with the biased 85% or the discredited 0.85 JSON.

\documentclass[conference]{IEEEtran}
\IEEEoverridecommandlockouts
\usepackage{cite}
\usepackage{amsmath,amssymb,amsfonts}
%% \usepackage{algorithmic}
\usepackage{graphicx}
\usepackage{textcomp}
\usepackage{xcolor}
\usepackage{booktabs}
\usepackage{array}
\def\BibTeX{{\rm B\kern-.05em{\sc i\kern-.025em b}\kern-.08em T\kern-.1667em\lower.7ex\hbox{E}\kern-.125emX}}

\begin{document}

\title{Does Fusing Deformation and Shear Improve Tactile Material
Classification? An Ablation Study under Contact-Condition Variability}

\author{\IEEEauthorblockN{Amr Magdy Mohamed Elsayed Abdalla}
\IEEEauthorblockA{\textit{DIBRIS, University of Genoa} \\
Cognitive Architectures for Robotics (C2b)\\
}
\and
\IEEEauthorblockN{PROFESSOR MASTROGIOVANNI FULVIO , Gianluca Galvagni}
\IEEEauthorblockA{\textit{DIBRIS, University of Genoa}\\
Genoa, Italy}
}

\maketitle

\begin{abstract}
Tactile sensing lets a robot identify what it touches from the mechanical
response of the contact, independent of vision. We study whether \emph{fusing}
the two channels of a vision-based tactile sensor---surface deformation and
shear---improves material classification over either channel alone, and how
robustly that advantage holds as contact conditions vary. Using a Daimon WS
sensor we collected 10 materials under two contact modes (a controlled
\emph{press} and a more variable \emph{airhold} squeeze), and compared classical
classifiers and deep temporal models under a leakage-free, validation-based
evaluation. Three findings emerge. First, the \emph{spatial} structure of the
contact footprint dominates raw force magnitude (65.5\% vs.\ 42.8\% accuracy),
confirming that material identity is encoded geometrically. Second, fusing
deformation and shear beats the best single modality by 4--9 percentage points
in every condition, because compliant and textured materials activate the shear
channel. Third, a temporal self-attention model raises honest held-out accuracy
from a $\sim$71\% classical ceiling to $\sim$80\%, while a controlled
nine-configuration regularization sweep shows the residual train--test gap is
driven by data quantity rather than model capacity. Visually identical PLA cubes
of differing stiffness are separated by touch with near-zero cross-confusion. We
report a reproducibility-bounded result of $\sim$80\%~$\pm$~5--6\% on 10
materials and discuss the hand-collection force confound as the principal
limitation.
\end{abstract}

\begin{IEEEkeywords}
tactile sensing, material classification, deformation, shear, sensor fusion,
temporal attention, robotic manipulation
\end{IEEEkeywords}

% =====================================================================
\section{Introduction and Research Question}
% =====================================================================
Touch gives a robot information vision cannot: the mechanical response of a
contact reveals stiffness, texture, and compliance even when two objects look
identical. Vision-based tactile sensors expose this response as two physically
distinct channels---\emph{deformation} (the normal/vertical displacement of a
soft sensing surface, encoding contact geometry) and \emph{shear} (the
lateral/tangential displacement, proportional to tangential force)
\cite{gelsight}. A natural question for any fusion-based perception system is
whether using both channels is worth the cost over using one.

Following the PICO formulation for robotics research questions, we phrase our
study as: \emph{Does fusing deformation and shear (intervention) improve
material-classification accuracy (outcome) on a set of everyday rigid, compliant,
and textured objects (population/context) compared to using either modality
alone (baseline), and how does that advantage change as contact-condition
variability increases?} This is testable, measurable (classification accuracy
under cross-validation), controlled (same sensor, materials, and pipeline across
conditions), and non-trivial (prior work has not run the controlled
single-vs-fused comparison on this modality).

We formalize three hypotheses:
\begin{itemize}
\item \textbf{H1}: fusion outperforms the best single modality at low contact
variability;
\item \textbf{H2}: the fused advantage persists, and is not eliminated, as
contact variability grows;
\item \textbf{H3}: each single modality fails on a physically predictable subset
of materials, which fusion covers.
\end{itemize}

\noindent\textbf{Variables.} The \emph{independent} variables are the sensing
modality (deformation-only / shear-only / fused), the feature representation
(magnitude / shape / spatial), the model class (classical / temporal deep), and
the contact condition (press / airhold). The \emph{dependent} variable is
5-fold cross-validated classification accuracy (with per-class confusion and the
train--test gap as secondary metrics). \emph{Control} variables are the sensor,
the material set, the trial count, and the evaluation protocol, held fixed
across comparisons.

\noindent\textbf{Contributions.} (i) A controlled feature-group and modality
ablation that isolates the contribution of spatial structure and of shear; (ii)
a demonstration that temporal self-attention---not added model
capacity---breaks the classical accuracy ceiling; (iii) an honest,
reproducibility-bounded evaluation with a regularization sweep that diagnoses
the residual generalization gap as data-limited; and (iv) a clear
``touch-beats-vision'' result on visually identical objects.

% =====================================================================
\section{Literature Analysis}
% =====================================================================
Our study sits on a ladder of four works spanning sensing principles to the
current state of the art.

\textbf{Foundation.} GelSight \cite{gelsight} established the vision-based
tactile paradigm: a camera images a soft elastomer whose vertical and lateral
surface deformation encodes contact geometry, with marker-motion magnitude
``roughly proportional to'' contact force and shear. It grounds our treatment of
deformation and shear as the two physically meaningful channels and motivates
modeling the temporal loading dynamics of a soft, viscoelastic contact.

\textbf{Direct lab/pipeline precedent.} Amin, Gianoglio and Valle
\cite{amin2023}, from the same institution on the same Baxter platform, perform
embedded real-time object \emph{hardness} classification, comparing an SVM,
single-layer feed-forward networks, and a CNN (CNN best, $>$98\%). Crucially
their sensor is a \emph{PVDF piezoelectric pressure array}---not a
deformation+shear field---and their features are the mean and standard deviation
of each tactile frame. Our work mirrors their Baxter-mounted classical-plus-deep
pipeline but differs in modality and, importantly, shows that the per-frame
mean/std pooling they adopt discards exactly the spatial structure that we find
carries material identity. (We note for accuracy that the ``170 trials'' in
\cite{amin2023} refer to materials-machine compression ground-truthing, not
tactile-dataset size.)

\textbf{Sensor-specific reference.} Liao et al.\ \cite{liao2025} use the Daimon
DM-Tac sensor---the same vendor family as ours---and decompose contact into
normal and shear components, estimating fruit hardness from the
\emph{dynamics of the normal force} ($\mathrm{d}F_z/\mathrm{d}t$,
$\mathrm{d}F_z/\mathrm{d}\delta$). This validates our sensor lineage and supports
modeling temporal loading dynamics rather than a single peak frame. Notably,
\cite{liao2025} \emph{decomposes shear but uses only the normal force},
motivating precisely the fusion question we test.

\textbf{State-of-the-art ceiling.} FG-CLTP \cite{fgcltp2025} is a fine-grained
contrastive language--tactile pretraining framework that aligns 3D tactile point
clouds with quantitative contact-state descriptions, reporting 95.9\% accuracy
on a 100k-pair dataset. Its central argument---that \emph{quantitative} contact
states (force magnitude, contact geometry) carry more than qualitative texture
descriptors, and that spatial 3D representations outperform image-only
ones---directly supports our premise. It also delimits the resource frontier
(data scale, 3D representation, pretraining) that our small single-sensor study
does not reach, framing our future work.

\textbf{Positioning.} Among these, ours is the only study that (a) uses
deformation+shear as its explicit fused modality, (b) runs the controlled
deformation/shear/fused ablation that none of the four perform, and (c) reports
under an honest, reproducibility-bounded protocol. Where \cite{amin2023} averages
the spatial field away and \cite{liao2025} discards shear, our contribution is to
quantify what is recovered by keeping both.

% =====================================================================
\section{Methodological Proposal}
% =====================================================================
\subsection{Sensor and signals}
We use a Daimon WS vision-based tactile sensor providing, per frame, a
deformation field and a shear field, each a two-component vector map downsampled
to $30\times40$. A trial is a temporal stack of such frames spanning approach,
loading, hold, and release. Concatenating the two modalities yields a
$(T,30,40,4)$ tensor per trial.

\subsection{Feature representations (classical)}
We extract a 30-dimensional feature set in four groups: \emph{magnitude}
(peak, mean, std of the field), \emph{shape} (loading-curve descriptors),
\emph{spatial-deformation} (centroid, spread, area, concentration of the
deformation footprint), and \emph{spatial-shear} (the same for shear). The
feature-group split lets us ablate which physical aspect of the contact carries
identity.

\subsection{Models}
\emph{Classical}: SVM (RBF), Random Forest, and XGBoost on the 30 features.
\emph{Deep, temporal}: a CNN encodes each frame; a bidirectional LSTM models the
loading sequence; a temporal self-attention layer re-weights informative frames
before a linear classifier. We additionally test a dual-branch variant that
processes deformation and shear in separate streams, to isolate whether explicit
modality separation helps.

\subsection{Honest evaluation protocol}
All accuracies are 5-fold cross-validated. For the deep models, early stopping
uses an \emph{internal validation split carved from the training fold}; the test
fold is never used for model selection. This removes the optimistic bias of
selecting the epoch by test-fold accuracy and yields reviewer-proof numbers. We
report mean accuracy, standard deviation across folds, the train--test gap, and
the aggregated confusion matrix.

% =====================================================================
\section{Experiment Design}
% =====================================================================
\subsection{Materials and contact conditions}
We collected 10 materials spanning rigid, compliant, and textured classes:
carton box, ceramic cup, hardwood, two visually identical PLA cubes of differing
stiffness (\texttt{pla\_hard}, \texttt{pla\_soft}), two textured plastics, rubber,
steel cup, and a wallet. Each material was sampled under two contact conditions,
which constitute our variability axis: \emph{press} (a controlled, consistent
contact) and \emph{airhold} (the object squeezed against the sensor in the air,
producing more variable contact and stronger shear). Twenty trials per material
per condition were collected.

\subsection{Collection constraint and its consequence}
Robot access was limited, so data were hand-collected with a release-gated
protocol guaranteeing one clean press per trial. This makes the applied force
\emph{uncontrolled}, so force is partially confounded with material identity.
This is the principal limitation of the study and is treated explicitly: results
are validation-grade rather than force-controlled-deployment-grade. The two
contact conditions are used precisely to probe robustness to this variability
(H2).

\subsection{Baseline}
Per the RT2 methodology, the baseline is the ``most common solution'': a
classical classifier on magnitude features (the mean/std-style representation of
\cite{amin2023}) and a standard CNN-LSTM without temporal attention. The
intervention (spatial features; modality fusion; temporal attention) is compared
against these.

% =====================================================================
\section{Results and Visualization}
% =====================================================================
\subsection{Spatial structure dominates magnitude}
A feature-group ablation (Table~\ref{tab:feat}) shows spatial features reach
65.5\% on the pooled set versus 42.8\% for magnitude---a $\sim$1.5$\times$
advantage, robust across conditions. Exploratory analysis
(Fig.~\ref{fig:confound}) explains why: raw loading curves separate materials by
force, but after per-trial normalization the curves partially collapse, leaving
residual \emph{shape} differences as the discriminative signal. Per-material
contact imprints (Fig.~\ref{fig:imprints}) are visibly distinct, confirming that
identity is encoded geometrically.

\begin{table}[t]
\caption{Feature-group ablation (5-fold CV accuracy, SVM).}
\label{tab:feat}
\centering
\begin{tabular}{lcccc}
\toprule
Group (dim) & press & airhold & all \\
\midrule
Magnitude (7) & 50.5 & 44.0 & 42.8 \\
Shape (5) & 27.0 & 33.5 & 26.2 \\
\textbf{Spatial (18)} & \textbf{68.5} & \textbf{67.0} & \textbf{65.5} \\
ALL (30) & 74.5 & 72.5 & 71.0 \\
\bottomrule
\end{tabular}
\end{table}

\subsection{Fusion beats single modality (answer to the research question)}
The modality ablation (Table~\ref{tab:mod}, Fig.~\ref{fig:modality}) confirms
H1: fusing deformation and shear beats the best single modality by 4--9 points
in every condition. H3 holds in the confusion structure: deformation struggles
where stiffness is similar, shear is weak on smooth rigid items, and fusion
covers both. Feature-importance analysis ranks the spatial spread of \emph{both}
fields highest, and the shear channel---negligible on rigid pilot materials---
becomes informative once compliant and textured materials are present
(supporting H2).

\begin{table}[t]
\caption{Modality ablation (5-fold CV accuracy, SVM).}
\label{tab:mod}
\centering
\begin{tabular}{lccc}
\toprule
Modality & press & airhold & all \\
\midrule
Deformation only & 67.5 & 68.0 & 62.0 \\
Shear only & 65.5 & 62.5 & 55.5 \\
\textbf{Fused} & \textbf{74.5} & \textbf{72.5} & \textbf{71.0} \\
\midrule
\emph{Gain over best single} & +7.0 & +4.5 & +9.0 \\
\bottomrule
\end{tabular}
\end{table}

\subsection{Temporal attention breaks the ceiling}
Classical models and a baseline CNN-LSTM converge near 71\%. Temporal
self-attention raises \emph{honest} held-out accuracy to $\sim$80\%
(press 80.5\%, airhold 79.5\%, pooled 79.2\%). A dual-branch variant matches but
does not exceed single-stream attention, isolating temporal attention---not
added capacity or explicit modality separation---as the source of the gain
(Fig.~\ref{fig:acc}).

\subsection{The gap is data-limited (regularization sweep)}
A nine-configuration sweep (Fig.~\ref{fig:reg}) finds that only gentle output
regularization helps: label smoothing gives the best single result (82.8\%),
weight decay helps marginally, while capacity reduction, temporal cropping, and
augmentation are neutral or harmful. The train--test gap stays at +12--17\%
under every intervention, and two runs of the identical architecture scored
82.8\% and 77.0\%---a $\sim$6-point swing from random seed alone. We therefore
report a reproducibility-bounded \textbf{$\sim$80\%~$\pm$~5--6\%} on 10
materials, with 82.8\% as the best single configuration. The immovable gap is
direct evidence that data quantity ($\sim$20 trials/class), not model capacity,
is the binding constraint.

\subsection{Touch beats vision}
The two PLA cubes are visually indistinguishable yet differ in stiffness. In the
honest confusion matrix (Fig.~\ref{fig:confusion}) they show near-zero
cross-confusion---direct evidence that the mechanical-response modality
discriminates what a camera cannot. Residual errors concentrate in
physically similar pairs (the rigid-flat cluster; the two cups), confirming the
classifier's mistakes track genuine material similarity.

% ---- figure stubs: replace with generated PNGs in reports_2/figures ----
\begin{figure}[t]\centering
\includegraphics[width=\columnwidth]{reports_2/figures/magnitude_confound_data2.png}
\caption{Raw loading curves separate materials by force; normalized curves
partially collapse, leaving shape as the discriminative signal.}
\label{fig:confound}\end{figure}

\begin{figure}[t]\centering
\includegraphics[width=\columnwidth]{reports_2/figures/spatial_imprints_data2.png}
\caption{Per-material contact imprints (onset/peak/release; press and airhold).
Each material leaves a distinct spatial footprint.}
\label{fig:imprints}\end{figure}

\begin{figure}[t]\centering
\includegraphics[width=\columnwidth]{reports_2/figures/modality_ablation.png}
\caption{Modality ablation: fused deformation+shear beats either alone in every
condition (annotated gain over best single modality).}
\label{fig:modality}\end{figure}

\begin{figure}[t]\centering
\includegraphics[width=\columnwidth]{reports_2/figures/accuracy_comparison.png}
\caption{Accuracy by model and condition under honest evaluation. Attention
breaks the $\sim$71\% classical ceiling; label smoothing is best.}
\label{fig:acc}\end{figure}

\begin{figure}[t]\centering
\includegraphics[width=\columnwidth]{reports_2/figures/regularization_sweep.png}
\caption{Regularization sweep (10 materials, honest eval). Only gentle output
regularization helps; the $\sim$80\%~$\pm$~5--6\% band is shaded.}
\label{fig:reg}\end{figure}

\begin{figure}[t]\centering
\includegraphics[width=\columnwidth]{reports_2/figures/confusion_matrix_honest.png}
\caption{Honest attention confusion matrix (pooled, 10 materials). The PLA pair
shows near-zero cross-confusion---``touch beats vision''.}
\label{fig:confusion}\end{figure}

% =====================================================================
\section{Discussion and Limitations}
% =====================================================================
\textbf{Answer to the research question.} Fusing deformation and shear improves
material classification over either modality alone by 4--9 points (H1), the
advantage persists across both contact conditions and is driven by compliant and
textured materials activating shear (H2, H3). The underlying signal is
\emph{spatial}: contact-footprint geometry, not raw force magnitude.

\textbf{What attention adds, and what it does not.} Temporal self-attention
recovers $\sim$9 points over the classical/baseline ceiling, showing the
limitation there was partly temporal-modeling, not purely data. But the
nine-configuration sweep shows the residual generalization gap is
\emph{data-quantity-driven}: it is immovable by capacity reduction or
augmentation, and only gentle output regularization helps marginally.

\textbf{Limitations.} (1) \emph{Hand collection} leaves applied force confounded
with material identity; this caps accuracy and is the dominant limitation.
(2) \emph{Small data} ($\sim$20 trials/class) yields wide fold variance
($\pm$5--6\%) and a $\sim$14\% train--test gap, motivating band-based reporting.
(3) \emph{Tactile-specific augmentation constraint}: because the label is encoded
in spatial geometry and magnitude, standard image augmentations (rotation,
magnitude scaling, even $\pm$1\,px translation) corrupt the label and degrade
accuracy---only temporal jitter and small sensor-scale noise are
label-preserving. (4) \emph{Offline, single-sensor, single-operator}: no
inter-session or on-robot deployment results yet.

\textbf{Future work.} Constant-force robotic collection (removing the confound),
more trials per class (closing the gap), an on-robot deployment/transfer test,
and richer 3D representations or contrastive pretraining in the direction of
\cite{fgcltp2025}.

\textbf{Conclusion.} On 10 materials with a deformation+shear tactile sensor,
fusing the two modalities and modeling the temporal loading dynamics with
attention yields $\sim$80\%~$\pm$~5--6\% honest accuracy, separating even
visually identical objects by touch. The result is a complete, self-aware answer
to the research question, bounded by data scale rather than by method.

% =====================================================================
\section*{Reproducibility}
Code, per-fold logs, and figure-generation scripts are in the project
repository. Figures are regenerated from the dataset and result files;
all reported numbers use the honest validation-based protocol.

\begin{thebibliography}{00}
\bibitem{gelsight} W.~Yuan, S.~Dong, and E.~H.~Adelson, ``GelSight:
High-Resolution Robot Tactile Sensors for Estimating Geometry and Force,''
\emph{Sensors}, vol.~17, no.~12, 2762, 2017.
\bibitem{amin2023} Y.~Amin, C.~Gianoglio, and M.~Valle, ``Embedded real-time
objects' hardness classification for robotic grippers,'' \emph{Future
Generation Computer Systems}, vol.~148, pp.~211--224, 2023.
\bibitem{liao2025} Z.~Liao, Y.~Du, J.~Duan, H.~Liang, and M.~Y.~Wang,
``Quantitative Hardness Assessment with Vision-based Tactile Sensing for Fruit
Classification and Grasping,'' 2025.
\bibitem{fgcltp2025} W.~Ma, C.~Zhang, Y.~Cai, G.~Yao, S.~Cui, and S.~Wang,
``FG-CLTP: Fine-Grained Contrastive Language Tactile Pretraining for Robotic
Manipulation,'' 2025.
\end{thebibliography}

\end{document}
